\documentclass[../../../include/open-logic-section]{subfiles}

\begin{document}

\olfileid{sth}{choice}{vitali}
\olsection{అనుబంధం: విటాలి వైరుధ్యాభాసం}

బనాక్--టార్స్కీ నిర్మాణం ఆమోదయోగ్యమా కాదా అనేది నిజంగా అర్థం కావాలంటే దాని \emph{నిరూపణ}ను పరిశీలించాలి. దురదృష్టవశాత్తూ అందుకు ఇక్కడ ఇవ్వగలిగినదానికంటే చాలా ఎక్కువ బీజగణితం కావాలి. అయితే తరువాతి ఫలితంపై దృష్టి పెట్టి, బనాక్--టార్స్కీ నిరూపణను అర్థం చేసుకోవడానికి తోడ్పడే కొన్ని సంక్షిప్త వ్యాఖ్యలు ఇవ్వవచ్చు.\footnote{మరింత విస్తృతమైన వివరణకు \cite{Weston2003} లేదా \cite{Wagon2016} చూడండి.}

\begin{thm}[విటాలి వైరుధ్యాభాసం ($\ZFC$లో)]\ollabel{vitaliparadox}
ఏ వృత్తాన్నైనా లెక్కించదగిన సంఖ్యలో ముక్కలుగా విడదీసి, ఆ ముక్కలను భ్రమణం, స్థానాంతరణం ద్వారా మళ్లీ అమర్చి, ఆ వృత్తానికి రెండు ప్రతులను ఏర్పరచవచ్చు.
\end{thm}

విటాలి వైరుధ్యాభాసాన్ని నిరూపించడం బనాక్--టార్స్కీ వైరుధ్యాభాసం కంటే సులభం. కొలవలేని సమితికి విటాలి ఇచ్చిన \citeyear{Vitali1905} నిర్మాణం నుంచి ఇది వస్తుంది కాబట్టి దీన్ని ``విటాలి వైరుధ్యాభాసం'' అన్నాం. అయితే విటాలి, బనాక్--టార్స్కీ నిరూపణలలో సమితి-సిద్ధాంత భాగాలు చాలా పోలి ఉంటాయి. ప్రధాన తేడా ఏమిటంటే బనాక్--టార్స్కీలో \emph{పరిమిత} విభజన, విటాలి వైరుధ్యాభాసంలో \emph{లెక్కించదగిన అనంత} విభజన. \citet{Weston2003} చెప్పినట్లుగా, విటాలి వైరుధ్యాభాసం ``బనాక్--టార్స్కీ అంత ఆశ్చర్యకరం కాదు; కానీ `సరళ' పరిస్థితుల్లోనూ జ్యామితీయ వైరుధ్యాభాసాలు రావచ్చని చూపుతుంది.''

విటాలి వైరుధ్యాభాసం ద్విమితీయ ఆకృతియైన వృత్తానికి సంబంధించినది. కాబట్టి $\Real^2$ సమతలంలో పనిచేద్దాం. $\rotationsgroup$ అనేది మూలబిందువు చుట్టూ సవ్యదిశలో, $[0,2\pi)$ పరిధిలో పూర్తి చుట్టుకు పరిమేయ గుణితాలైన రేడియన్ కోణాలతో చేసే బిందువుల భ్రమణాల సమితి. $\rotationsgroup$ గురించి కొన్ని బీజగణిత విషయాలు ఇవి (వాక్యం అర్థం కాకపోతే నిరూపణ ద్వారా స్పష్టమవుతుంది):
\sourcecorrection{OLTESTCHOICEVITALI-001}{మూలం $[0,2\pi)$లో పరిమేయ రేడియన్ విలువలనే తీసుకుంది; అవి $2\pi$ మాడ్యులో భ్రమణ సమూహంగా మూసుకోకపోవచ్చు, ప్రతిలోమం కూడా ఉండకపోవచ్చు. కోణాలను $2\pi$కు పరిమేయ గుణితాలుగా సరిచేశాం; వృత్తంలోని పరిమేయ-మలుపుల సమూహమే తరువాతి వాదనకు అవసరం.}

\begin{lem}\ollabel{rotationsgroupabelian}
ప్రమేయాల సంయోజన క్రింద $\rotationsgroup$ ఒక అబేలియన్ \emph{సమూహం}.
\end{lem}

\begin{proof}
$0$ రేడియన్ల భ్రమణాన్ని $0_{\rotationsgroup}$గా రాస్తే, అది $\rotationsgroup$కు తటస్థ మూలకం; ఎందుకంటే ఏ $\rho \in \rotationsgroup$కైనా $\comp{0_{\rotationsgroup}}{\rho} = \comp{\rho}{0_{\rotationsgroup}} =
\rho$.

ప్రతి మూలకానికీ ప్రతిలోమం ఉంది. $\rho \in \rotationsgroup$ భ్రమణం $r$ రేడియన్లు అయితే, శూన్య భ్రమణం సందర్భంలో దాని ప్రతిలోమం తానే; లేకపోతే $\rho^{-1} \in \rotationsgroup$ అనేది $2\pi - r$ రేడియన్ల భ్రమణం. అందువల్ల $\rho \circ \rho^{-1} = 0_\rotationsgroup$.
\sourcecorrection{OLTESTCHOICEVITALI-002}{మూలం శూన్య భ్రమణానికి ప్రతిలోమాన్ని $2\pi-r=2\pi$గా రాసింది; అది ఎంచుకున్న $[0,2\pi)$ పరిధిలో లేదు. $r=0$ను విడిగా చూపి, మిగిలిన సందర్భానికి $2\pi-r$ను వాడాం.}

సంయోజనకు సహచర్య ధర్మం ఉంది: ఏ $\rho, \sigma, \tau \in
\rotationsgroup$కైనా $\comp{\rho}{(\comp{\sigma}{\tau})} =
\comp{(\comp{\rho}{\sigma})}{\tau}$.

సంయోజనకు క్రమమార్పిడి ధర్మం ఉంది: ఏ $\rho, \sigma \in \rotationsgroup$కైనా $\comp{\rho}{\sigma} =
\comp{\sigma}{\rho}$.
\end{proof}

నిజానికి మన సమూహం $\rotationsgroup$ను రెండు భాగాలుగా విడదీసి, వాటిలో ఏ భాగం నుంచైనా మొత్తం సమూహాన్ని తిరిగి పొందవచ్చు:

\begin{lem}\ollabel{disjointgroup}
$\rotationsgroup$ను $\rotationsgroup_{1}$, $\rotationsgroup_{2}$ అనే రెండు విచ్ఛిన్న సమితులుగా విభజించవచ్చు; రెండూ $\rotationsgroup$కు జనక సమితులుగా పనిచేస్తాయి.
\end{lem}

\begin{proof}
$\rotationsgroup_{1}$లో $[0, \pi)$ పరిధిలో పూర్తి చుట్టుకు పరిమేయ గుణితాలైన రేడియన్ కోణాల భ్రమణాలు ఉండాలి; $\rotationsgroup_{2} = \rotationsgroup
\setminus \rotationsgroup_1$. ప్రాథమిక బీజగణితం ప్రకారం $\Setabs{\comp{\rho}{\rho}}{\rho \in \rotationsgroup_1} =
\rotationsgroup$. $\rotationsgroup_2$కూ ఇలాంటి ఫలితం వస్తుంది.
\end{proof}

సమూహాల గురించిన ఈ విషయంతో \olref{vitaliparadox}ను నిరూపిద్దాం. $\onesphere$ ఏకక వృత్తం, అంటే సమతల మూలబిందువు నుంచి ఖచ్చితంగా~$1$ దూరంలో ఉన్న బిందువుల సమితి, అంటే $\Setabs{\tuple{r,s} \in \Real^2}{\sqrt{r^2+s^2}=1}$. $\onesphere$పై ఈ సంబంధాన్ని తీసుకొని $\onesphere$ను భాగాలుగా విభజిస్తాం:
\[ 
	r \sim s \emph{iff}(\exists \rho \in \rotationsgroup)\rho(r) = s.
\]
అంటే మూలబిందువు చుట్టూ పరిమేయ-మలుపు భ్రమణంతో ఒక బిందువు నుంచి మరొకదానికి వెళ్లగలిగితే, అప్పుడే $\onesphere$లోని ఆ రెండు బిందువులకు ఈ సంబంధం ఉంది. ఆశ్చర్యం లేకుండా:

\begin{lem}
$\sim$ సమానత్వ సంబంధం.
\end{lem}

\begin{proof}
\olref{rotationsgroupabelian}ను వాడితే నేరుగా వస్తుంది.
\end{proof}

ఇప్పుడు ఎంపికను ఉపయోగించి, $\onesphere$పై $\sim$ సంబంధంలోని ప్రతి సమానత్వ తరగతి నుంచి ఖచ్చితంగా ఒక్క మూలకాన్ని కలిగిన $C$ సమితిని పొందుతాం. అంటే సమానత్వ తరగతుల సమితిపై $f$ అనే ఎంపిక ప్రమేయాన్ని తీసుకుంటాం;\footnote{$\rotationsgroup$ \tetoken{లెక్కించదగినది}{!!{enumerable}} కాబట్టి $E$లోని ప్రతి \tetoken{మూలకం}{!!{element}} \tetoken{లెక్కించదగినది}{!!{enumerable}}. $\onesphere$ \tetoken{లెక్కించలేనిది}{!!{nonenumerable}} కాబట్టి \olref{vitalicover}, \olref[card-arithmetic][simp]{kappaunionkappasize} ప్రకారం $E$ కూడా \tetoken{లెక్కించలేనిది}{!!{nonenumerable}}. కాబట్టి ఇది లెక్కించదగిన ఎంపికను మించిన ఎంపిక వినియోగం.}
\[
	E = \Setabs{\equivrep{r}{\sim}}{r \in \onesphere},
\]
అప్పుడు $C = \ran{f}$. ప్రతి భ్రమణం $\rho \in \rotationsgroup$కు $\funimage{\rho}{C}$ అనేది $C$లోని ప్రతి బిందువుపై $\rho$ను వర్తింపజేసి పొందిన బిందువుల సమితి. తరువాతి రెండు ఫలితాలు ఈ సమితులు వృత్తాన్ని పూర్తిగా, ఒకదానితో మరొకటి అతివ్యాప్తి లేకుండా కప్పుతాయని చూపుతాయి:

\begin{lem}\ollabel{vitalicover}
$\onesphere = \bigcup_{\rho \in \rotationsgroup} \funimage{\rho}{C}$.
\end{lem}

\begin{proof}
$s \in \onesphere$ను స్థిరపరచుకుందాం. $r \in
\equivrep{s}{\sim}$ అయ్యే ఏదో $r \in C$ ఉంది; అంటే $r \sim s$; అంటే ఏదో $\rho \in \rotationsgroup$కు $\rho(r) = s$.
\end{proof}

\begin{lem}\ollabel{vitalinooverlap}
$\rho_1 \neq \rho_2$ అయితే $\funimage{\rho_1}{C} \cap \funimage{\rho_2}{C} = \emptyset$.
\end{lem}

\begin{proof}
$s \in \funimage{\rho_{1}}{C} \cap \funimage{\rho_{2}}{C}$ అనుకుందాం. అప్పుడు ఏవో $r_{1}, r_{2} \in C$కు $s = \rho_{1}(r_{1}) = \rho_{2}(r_{2})$. అందువల్ల $\rho^{-1}_2(\rho_1(r_1)) = r_2$; ఇంకా $\comp{\rho_1}{\rho^{-1}_2} \in \rotationsgroup$, కాబట్టి $r_{1} \sim
r_{2}$. $C$ అనేది $\sim$ క్రింద ప్రతి సమానత్వ తరగతి నుంచి ఒక్క మూలకాన్నే ఎంచుకుంటుంది కనుక $r_1 = r_2$. అప్పుడు $s = \rho_1(r_1) = \rho_2(r_1)$; వృత్తంపై వేర్వేరు భ్రమణాలు ఒక్క బిందువునూ ఒకే స్థానానికి పంపవు కాబట్టి $\rho_1 = \rho_2$.
\end{proof}

ఇప్పుడు ముందరి బీజగణిత విషయాలను వృత్తానికి వర్తింపజేద్దాం:

\begin{lem}\ollabel{pseudobanachtarski}
$\onesphere$ను $D_{1}$, $D_{2}$ అనే రెండు విచ్ఛిన్న సమితులుగా విభజించవచ్చు; $D_{1}$ను లెక్కించదగిన సంఖ్యలో సమితులుగా విభజించి, వాటిని భ్రమణం ద్వారా $\onesphere$ ప్రతిగా అమర్చవచ్చు ($D_{2}$కు కూడా అలాగే).
\end{lem}

\begin{proof}
\olref{disjointgroup}లోని $\rotationsgroup_{1}$, $\rotationsgroup_{2}$లను వాడి ఇలా తీసుకుందాం:
\begin{align*}
	D_{1} &= \bigcup_{\rho \in \rotationsgroup_1} \funimage{\rho}{C} &
	D_{2} &= \bigcup_{\rho \in \rotationsgroup_2} \funimage{\rho}{C}
\end{align*}
\olref{vitalicover} ప్రకారం ఇది $\onesphere$ విభజన; \olref{vitalinooverlap} ప్రకారం $D_1$, $D_2$ విచ్ఛిన్నాలు. నిర్మాణం ప్రకారం ప్రతి $\rho \in \rotationsgroup_1$కు $\funimage{\rho}{C}$గా $D_1$ను లెక్కించదగిన సంఖ్యలో సమితులుగా విభజించవచ్చు. వీటిని భ్రమణం ద్వారా $\onesphere$ ప్రతిగా అమర్చవచ్చు: ఎందుకంటే \olref{disjointgroup}, \olref{vitalicover} ప్రకారం $\onesphere = \bigcup_{\rho \in
\rotationsgroup}\funimage{\rho}{C} = \bigcup_{\rho \in
\rotationsgroup_1}\funimage{(\comp{\rho}{\rho})}{C}$. $D_2$కూ అదే వాదన వర్తిస్తుంది. \sourcecorrection{OLTESTCHOICEVITALI-003}{మూల గద్యంలో $\rho\in R_1$ అని ఉంది; ఈ నిర్మాణంలో నిర్వచించిన సమితి $\rotationsgroup_1$ మాత్రమే. దానిని సరైన సమూహ భాగంగా సరిచేశాం.}\end{proof}\noindent దీని నుంచి విటాలి వైరుధ్యాభాసం వెంటనే వస్తుంది. $\onesphere$ను లెక్కించదగిన సంఖ్యలో ముక్కలుగా (వివిధ $\funimage{\rho}{C}$లుగా) విభజించి, వాటిని దృఢ చలనాలతో కదిలిస్తే $\onesphere$కు \emph{రెండు} ప్రతులు వస్తాయి: $D_1$లోని ప్రతి ముక్కను భ్రమణం చేయాలి; $D_2$లోని ప్రతి ముక్కను ముందుగా స్థానాంతరణం చేసి, తరువాత భ్రమణం చేయాలి.

నిరూపణ పద్ధతిని మళ్లీ చూద్దాం. సమతలంలో భ్రమణాల సమూహం గురించి కొన్ని బీజగణిత విషయాలతో మొదలెట్టాం. ఆ సమూహాన్ని వాడి $\onesphere$ను సమానత్వ తరగతులుగా విభజించాం. తరువాత ప్రతి తరగతి నుంచి మూలకాన్ని ఎంపికతో ఎంచుకొని ``వైరుధ్యాభాసం''కు చేరుకున్నాం.

బనాక్--టార్స్కీని నిరూపించడానికి కూడా ఇదే పద్ధతిని వాడతాం. ప్రధాన తేడా: దానికవసరమైన బీజగణిత విషయాలు విటాలి వైరుధ్యాభాసం కంటే చాలా క్లిష్టమైనవి. అయితే ఆ బీజగణిత విషయాలకు ఎంపికతో సంబంధం లేదు. వాటిని సంక్షిప్తంగా చెబుతాం.

బనాక్--టార్స్కీ కోసం మొదట \olref{disjointgroup}కు సమానమైన విషయాన్ని స్థాపిస్తాం: కనీసం రెండు జనకాలు గల \emph{స్వేచ్ఛా సమూహాన్ని} నాలుగు భాగాలుగా విడదీసి, సహజంగా వాటిని ``కదిలించి'' మొత్తం సమూహానికి రెండు ప్రతులను పొందవచ్చు.\footnote{\emph{నాలుగు} భాగాలు చాలనే విషయం \cite{Robinson1947} వల్ల వచ్చింది. ఇటీవలి నిరూపణకు \citet[Theorem 5.2]{Wagon2016} చూడండి. సమూహపు భాగాలను ``కదిలించడం'' అని చెప్పడంలో \citet[p.~3]{Weston2003}ను అనుసరిస్తున్నాం.} తరువాత $\Real^3$ మూలబిందువు చుట్టూ రెండు ప్రత్యేక భ్రమణాల ద్వారా భ్రమణాల స్వేచ్ఛా సమూహం $F$ను జనింపజేయవచ్చని చూపుతాం.\footnote{\citet[Theorem 2.1]{Wagon2016} చూడండి.} (ఇక్కడివరకు ఎంపిక వాడలేదు.) ఇప్పుడు గోళం ఉపరితలంలోని బిందువులను, $F$లోని భ్రమణం ద్వారా ఒకటి నుంచి మరొకటి పొందగలిగితే, ``సమానమైనవి''గా భావిస్తాం. తరువాత ``సమానమైన'' బిందువుల ప్రతి సమానత్వ తరగతి నుంచి ఖచ్చితంగా ఒక బిందువును \emph{ఎంపికతో} ఎంచుకుంటాం. ఇక్కడ భ్రమణాలకు స్థిరంగా ఉండే అపవాద బిందువులను విడిగా పరిష్కరించాల్సి ఉంటుంది; ఆ దశను ఈ సంక్షిప్త రూపరేఖ ఇవ్వదు. ఆ మినహాయింపును పరిష్కరించిన తరువాత, \olref{pseudobanachtarski}లో వలె $F$ విభజనను గోళ ఉపరితలానికి వర్తింపజేసి, ఉపరితలాన్ని నాలుగు భాగాలుగా విడదీసి, వాటిని ``కదిలించి'' ఉపరితలానికి రెండు ప్రతులను పొందుతాం. ఇది \citep{Hausdorff1914}ను స్థాపిస్తుంది:
\sourcecorrection{OLTESTCHOICEVITALI-004}{మూలం ``ఏ స్వేచ్ఛా సమూహం'' అన్నా నాలుగు భాగాల వైరుధ్యాభాస విభజన సాధ్యమని చెప్పింది; ఒక జనకంతో స్వేచ్ఛా సమూహం అనంత చక్రీయ సమూహం, దీనికి ఆ వాదన వర్తించదు. కనీసం రెండు జనకాలు కావాలని అర్హత చేర్చాం. అంతేకాక గోళ ఉపరితలంపై భ్రమణ చర్యలో స్థిర బిందువుల అపవాద సమితిని మూల రూపరేఖ వదిలింది; దాన్ని విడిగా పరిష్కరించాల్సిన నిరూపణ ఖాళీగా ప్రకటించాం.}

\begin{thm}[హౌస్‌డార్ఫ్ వైరుధ్యాభాసం ($\ZFC$లో)]
ఏ గోళపు ఉపరితలాన్నైనా పరిమిత సంఖ్యలో ముక్కలుగా విడదీసి, భ్రమణం, స్థానాంతరణం ద్వారా వాటిని మళ్లీ అమర్చి, ఆ గోళపు ఉపరితలానికి రెండు విచ్ఛిన్న ప్రతులను ఏర్పరచవచ్చు.
\end{thm}

పూర్తి బనాక్--టార్స్కీ సిద్ధాంతం కోసం (అది గోళపు ఉపరితలానికే కాక లోపలికి కూడా సంబంధించినది) ఇంకొన్ని బీజగణిత యుక్తులు కావాలి. నిజం చెప్పాలంటే, బీజగణితంలోని ప్రధాన భాగంపై ఇవి అలంకరణ మాత్రమే. అందుకే వెస్టన్ ఇలా రాస్తారు:
\begin{quote}
  [\ldots] స్వేచ్ఛా సమూహాల గురించిన ఫలితం బనాక్--టార్స్కీ వైరుధ్యాభాసం నిరూపణలో \emph{కీలక దశ}. ఈ దృష్టిలో బనాక్--టార్స్కీ వైరుధ్యాభాసం $\Real^3$ గురించి చెప్పేదానికంటే, [$\Real^3$లో స్థానాంతరణాలు, భ్రమణాల] సమూహపు సంక్లిష్టత గురించి ఎక్కువగా చెబుతుంది. \cite[p.~16]{Weston2003}
\end{quote}
అంటే \emph{పరిమిత} విభజన (బనాక్--టార్స్కీలో వలె) ఇవ్వగలమా, లేక \emph{లెక్కించదగిన అనంత} విభజనే (విటాలి వైరుధ్యాభాసంలో వలె) ఇవ్వగలమా అనేది ద్విమితీయ లేదా త్రిమితీయ పరిస్థితిలోని కొన్ని సమూహ-సిద్ధాంత విషయాలపై ఆధారపడుతుంది.

ఒప్పుకోవాలి: ఈ చివరి విషయం \olref[banach]{sec} చివర్లోని చమత్కారాన్ని కొంచెం పాడుచేస్తుంది. అది ద్విమితీయంగా రాసినది కాబట్టి ``Banach-Tarski''ని ``Banach-Tarski Banach-Tarski''గా మార్చాలంటే లెక్కించదగిన \emph{అనంత} సంఖ్యలో ముక్కలుగా విడదీయాలి. చమత్కారాన్ని సరిచేయాలంటే త్రిమితీయంగా రాయాలి. దాన్ని పాఠకుల అభ్యాసానికి వదిలేస్తాం.

చివరిగా ఒక వ్యాఖ్య. \olref[banach]{sec}లో గోళం నుంచి వచ్చే ``ముక్కలు'' \emph{కొలవదగినవి} కావని, చిత్రించలేని ``అనంతంగా చెదరిన'' సమూహాలై ఉంటాయని అన్నాం. \olref{pseudobanachtarski}లో ఎంపిక వాడినప్పుడు కూడా ఇదే నిజం. దీన్ని వివరించడం ఉపయోగకరం.

మళ్లీ కొంత నేపథ్యాన్ని రూపరేఖగా చెబుదాం (ఇది \emph{కేవలం} రూపరేఖే; \emph{కొలత} గురించి పాఠ్యపుస్తక వివరణ చదవవచ్చు). $X$ సమితిపై కొలతను నిర్వచించడమంటే, $X$పై ఏదో ``$\sigma$-బీజగణితం''లోని ప్రతి $E$కు $\mu(E) \in
\Real$ అనే విలువను ఇవ్వడం. వివరాలు ఇక్కడ అవసరం లేదు; అయితే కొలత విలువలు ఋణేతరాలని, $\mu$ లెక్కించదగిన సంకలనీయతను పాటించాలని గుర్తుంచుకోవాలి: విచ్ఛిన్న సమితుల లెక్కించదగిన సమ్మేళనం కొలత, వాటి విడివిడి కొలతల మొత్తం. అంటే $X_n$లు విచ్ఛిన్నాలైతే $\mu(\bigcup_{n < \omega} X_n) = \sum_{n < \omega}\mu(X_n)$. సమితి ``కొలవలేనిది'' అంటే దానికి తగినట్లు కొలత కేటాయించలేమని అర్థం. ఇప్పుడు ముందరి $\rotationsgroup$ను వాడి:

\begin{cor}[విటాలి]
$\mu(\onesphere) = 1$ అయ్యే, అలాగే $X$, $Y$ సమరూప ఆకృతులైతే $\mu(X) = \mu(Y)$ అయ్యే కొలత $\mu$ అనుకుందాం. అప్పుడు ప్రతి $\rho \in
\rotationsgroup$కూ $\funimage{\rho}{C}$ కొలవలేనిది.
\end{cor}

\begin{proof}
వైరుధ్యం కోసం కాదనుకుందాం. ఏదో $\sigma \in \rotationsgroup$, $r \in \Real$లకు $\mu(\funimage{\sigma}{C}) =
r$ అని తీసుకుందాం. ప్రతి $\rho \in \rotationsgroup$కు $\funimage{\rho}{C}$, $\funimage{\sigma}{C}$ సమరూప ఆకృతులు; కనుక ప్రతి $\rho \in
\rotationsgroup$కు $\mu(\funimage{\rho}{C}) = r$. \olref{vitalicover}, \olref{vitalinooverlap} ప్రకారం $\onesphere =
\bigcup_{\rho \in \rotationsgroup}\funimage{\rho}{C}$ అనేది జతలవారీ విచ్ఛిన్న సమితుల లెక్కించదగిన సమ్మేళనం. కాబట్టి లెక్కించదగిన సంకలనీయత ప్రకారం $\mu(\onesphere) = 1$ అనేది ప్రతి $\funimage{\rho}{C}$ కొలతల మొత్తం:
\sourcecorrection{OLTESTCHOICEVITALI-005}{మూల కొలత నిరూపణలో భ్రమణ సూచికను రెండుసార్లు $\rho\in C$గా వ్రాసింది; $C$ వృత్త బిందువుల సమితి, భ్రమణాలది కాదు. రెండు చోట్లా $\rho\in\rotationsgroup$గా సరిచేశాం.}
\[
	1 = \mu(\onesphere) = \sum_{\rho \in \rotationsgroup}\mu(\funimage{\rho}{C}) = \sum_{\rho \in \rotationsgroup}r
\]
కానీ $r = 0$ అయితే $\sum_{\rho \in \rotationsgroup}r = 0$; $r
> 0$ అయితే $\sum_{\rho \in \rotationsgroup}r = \infty$.
\end{proof}

\end{document}
